\documentclass{article}
\usepackage{amsmath, amssymb, amsthm}
\usepackage{hyperref}
\usepackage{geometry}
\geometry{margin=1in}

\title{Erd\H{o}s Problem \#130}
\date{Last updated: 2025-08-31}
\author{Source: erdosproblems.com}

\begin{document}
\maketitle

\noindent\textbf{Status:} OPEN \quad \textbf{No prize}

\noindent\textbf{Tags:} graph theory, chromatic number

\medskip

\noindent\textbf{Problem Statement:}

Let $A\subset\mathbb{R}^2$ be an infinite set which contains no three points on a line and no four points on a circle. Consider the graph with vertices the points in $A$, where two vertices are joined by an edge if and only if they are an integer distance apart. How large can the chromatic number and clique number of this graph be? In particular, can the chromatic number be infinite?




Asked by Andr\'{a}sfai and Erd\H{o}s. Erd\H{o}s \cite{Er97b} also asked where such a graph could contain an infinite complete graph, but this is impossible by an earlier result of Anning and Erd\H{o}s \cite{AnEr45}. See also [213] .




References

[AnEr45] Anning, Norman H. and Erd\H{o}s, Paul, Integral distances . Bull. Amer. Math. Soc. (1945), 598-600.

[Er97b] Erd\H{o}s, Paul, Some old and new problems in various branches of combinatorics . Discrete Math. (1997), 227-231.


\bigskip
\noindent\small{Source: \url{https://www.erdosproblems.com/130}}
\end{document}
